\section{Boundaries of the exact-arithmetic hypotheses}
\label{app:boundaries}

The positive theorems of this paper rest on three hypotheses that are easy to
weaken by accident: a supplied global curvature bound, exact rational input,
and the right notion of a common kernel when certificates are reduced. This
appendix records what happens near each boundary. Section~\ref{app:bnd-degenerate}
shows that a fully degenerate convex quartic minimum is rational and easy,
while partial degeneracy defeats the Newton refinement used for
Theorem~\ref{thm:exact-upper}; the general degenerate case remains open.
Section~\ref{app:bnd-nullvector} shows that restricting a semidefinite
feasibility problem along a nullvector of one selected feasible matrix can
destroy all rational feasible points, and that a common nullvector avoids
this. Section~\ref{app:bnd-uncertain} shows that, at singular equality
constraints, enclosures of function values can fail to justify any accurate
upper bound on the optimal value, even where a root exists robustly under
perturbation; exact input removes this obstruction.
Section~\ref{app:bnd-penalty} gives an exact penalty representation of
bounded quadratic problems by a polynomial-size quadratic lift and explains
why it lies outside the convex theorems. None of these results is a hardness
theorem.

\subsection{Degenerate convex quartics}
\label{app:bnd-degenerate}

The upper bound of Theorem~\ref{thm:exact-upper} uses a supplied global lower
bound on the Hessian. It enters the quadratic convergence region of Newton's
method with polynomially many printed bits and then reaches doubly
exponential accuracy by polynomially many Newton steps written as a short
circuit. Without a curvature bound, the optimal value of a globally convex
rational quartic on a nonempty rational polyhedron, when the optimal value
is finite, can still be approximated in polynomial time with an additive objective-gap
guarantee~\cite[Corollary~1.2]{SlotSteurerWiedmer2025}. That guarantee
concerns the value, not the distance to a minimizer, and the question here
is exact comparison. Without the curvature bound, one case is elementary.

\begin{proposition}[A fully degenerate minimum]
\label{prop:bnd-zero-hessian}
Let $f\in\Q[x_1,\ldots,x_n]$ have degree at most four and be convex on $\R^n$.
Suppose a global minimizer $p$ satisfies $\nabla^2f(p)=0$. Then
\[
\argmin f=\{x:D^3f(x)=0\},
\]
a nonempty rational affine space. A rational minimizer, the exact minimum
value, and a rational affine description of all minimizers are computed by
rational linear algebra in polynomial time; a unique minimizer is rational.
\end{proposition}

\begin{proof}
For fixed $u,v$, the polynomial $\tau\mapsto u^{\mathsf T}\nabla^2f(p+\tau
v)u$ has degree at most two, is nonnegative by convexity, and vanishes at
$\tau=0$; hence its linear coefficient $D^3f(p)[u,u,v]$ vanishes. By
polarization in $u$, $D^3f(p)=0$. Stationarity and the exact Taylor expansion
give $f(p+z)=f(p)+\mathcal Q(z)$ with $\mathcal Q(z)=\tfrac1{24}D^4f[z,z,z,z]$,
a nonnegative convex form of degree four, possibly zero. Since $D^3f$ is
affine and vanishes at $p$, its zero set is $p+\mathcal V$ with
$\mathcal V=\{v:D^4f[v,\cdot,\cdot,\cdot]=0\}$, and $\mathcal Q$ vanishes on
$\mathcal V$, so $p+\mathcal V$ consists of minimizers.

Conversely, let $\mathcal Q(v)=0$. Then $\mathcal Q(\tau v)=0$ for all $\tau$,
and for $0<\eta<1$ convexity and homogeneity give
\[
\mathcal Q(z+\tau v)\le(1-\eta)\mathcal Q\Bigl(\frac z{1-\eta}\Bigr)
+\eta\mathcal Q\Bigl(\frac{\tau v}\eta\Bigr)=(1-\eta)^{-3}\mathcal Q(z).
\]
Letting $\eta\downarrow0$ and exchanging the roles of $z$ and $z+\tau v$ gives
$\mathcal Q(z+\tau v)=\mathcal Q(z)$ for all $z,\tau$. Differentiating in
$\tau$ gives $D^4f[v,z,z,z]=0$ for all $z$, and polarization gives
$v\in\mathcal V$. The equations $D^3f(x)=0$ are at most $\binom{n+2}3$ rational
affine equations with coefficients of polynomial length, so Gaussian
elimination gives the claimed description, and evaluating $f$ at a rational
solution gives the minimum.
\end{proof}

The vanishing of $D^3f(p)$ also follows from a derivative inequality for
convex quartics~\cite{Nesterov2025}. Among globally convex inputs, the
hypothesis can be recognized: solve
$D^3f(x)=0$, pick a rational solution $q$, and test $\nabla f(q)=0$ and
$\nabla^2f(q)=0$. Global convexity itself is not tested. The proposition does
not apply when only some Hessian eigenvalues vanish, as in $g(x)+y^4$ with
an irrational minimizer of $g$.

Partial degeneracy breaks the refinement step. For $f(x)=x^4$, Newton's map
is $N(x)=2x/3$, so reaching error $2^{-2^s}$ from $x_0=1$ takes more than
$2^s/\log_2(3/2)$ steps. This only concerns the iteration count: the minimizer
is the rational number zero.

\begin{proposition}[A discontinuous Newton map]
\label{prop:bnd-newton}
The rational quartic
\[
f(x,y)=(x^2+y^2)^2+xy^2+y^2=x^4+y^4+(2x^2+x+1)y^2
\]
is convex on $\R^2$, has the unique minimizer $(0,0)$ with Hessian
$\diag(0,2)$ there, and is strongly convex near every other point. Its Newton
map $N$ satisfies $N(0,\tau)\to(1/2,0)$ as $\tau\to0$. The extrapolated map
$\mathcal T=3N\circ N-2N$, which removes the linear error $N(z)=2z/3$ of
homogeneous quartics, satisfies $\mathcal T_y(0,\tau)=\tfrac58\tau+O(\tau^3)$.
\end{proposition}

\begin{proof}
Since $2x^2+x+1>0$, the origin is the unique zero and minimizer. The Hessian
is
\begin{gather*}
\nabla^2f=\begin{pmatrix}12x^2+4y^2&(8x+2)y\\(8x+2)y&4x^2+12y^2+2x+2\end{pmatrix},\\
\det\nabla^2f=24x^2(2x^2+x+1)+4y^2(24x^2-6x+1)+48y^4,
\end{gather*}
and both quadratic factors are positive (discriminants $-7$ and $-60$). Away
from the origin the first diagonal entry and the determinant are positive, so
the Hessian is positive definite there, and $f$ is convex. Direct inversion
at $(0,\tau)$, $\tau\ne0$, gives
\[
N(0,\tau)=\Bigl(\frac{1-2\tau^2}{2(1+12\tau^2)},\
\frac{\tau(16\tau^2-1)}{2(1+12\tau^2)}\Bigr)
=\bigl(\tfrac12-7\tau^2+O(\tau^4),\ -\tfrac12\tau+O(\tau^3)\bigr).
\]
Near $(1/2,0)$, $N$ is analytic and $N_y$ is odd in $y$, because $f$ is even
in $y$. Expanding $N_y=y-[(\nabla^2f)^{-1}\nabla f]_y$ to first order in $y$
gives
\[
N_y(x,y)=h(x)y+O(y^3),\qquad h(x)=\frac{x(8x+2)}{6(2x^2+x+1)},\qquad
h(\tfrac12)=\tfrac14 .
\]
Hence $\mathcal T_y(0,\tau)=3h(\tfrac12)(-\tfrac12\tau)-2(-\tfrac12\tau)+
O(\tau^3)=\tfrac58\tau+O(\tau^3)$.
\end{proof}

Thus points arbitrarily close to the minimizer are sent to distance about
$1/2$, although the Hessian is positive definite at every nonoptimal point,
and the natural correction for homogeneous quartics does not converge
quadratically. The differentiability of $N$ at the solution that such a
correction assumes is absent. A third shortcut would change coordinates so
that the Hessian kernel at the minimizer becomes a coordinate subspace; it
fails over $\Q$.

\begin{proposition}[An irrational Hessian kernel]
\label{prop:bnd-kernel}
Let $F(t,u,v)=t^4+2t^2-4t+(u-tv)^2+(u^2+v^2)^2$ and let $a$ be the real root
of $a^3+a-1=0$. Then $F$ is convex on $\R^3$, its unique minimizer is
$(a,0,0)$, and the kernel of $\nabla^2F(a,0,0)$ is the line spanned by
$(0,a,1)$, which contains no nonzero rational vector.
\end{proposition}

\begin{proof}
The univariate part has derivative $4(t^3+t-1)$ and second derivative
$12t^2+4>0$, the remaining terms are nonnegative, and the last one vanishes
only at $u=v=0$; this gives the minimizer. Write $y=(u,v)$, $r=\norm y$ and
$\mathfrak b=(1,-t)$. The Hessian blocks are $H_{tt}=12t^2+4+2v^2$,
$H_{ty}=(-2v,\,-2u+4tv)$ and $H_{yy}=2\mathfrak b\mathfrak b^{\mathsf T}
+4r^2I+8yy^{\mathsf T}$. If $r>0$, then $H_{yy}\succeq4r^2I$ and
$\norm{H_{ty}}\le(2+4\abs t)r$, so the Schur complement is at least
$12t^2+4+2v^2-(1+2\abs t)^2\ge4t^2+2+2v^2>0$. If $r=0$, the cross block
vanishes and both diagonal blocks are positive semidefinite. At the
minimizer, $H_{ty}=0$, $H_{tt}=12a^2+4>0$, and $H_{yy}=2\mathfrak b\mathfrak
b^{\mathsf T}$ with $\mathfrak b=(1,-a)$, whose kernel is spanned by $(a,1)$.
The cubic has no rational root, so $a$ is irrational, and a rational multiple
of $(0,a,1)$ with nonzero third coordinate would make $a$ rational.
\end{proof}

Approximate or algebraic coordinates adapted to the kernel are not excluded.

Quadratic regularization gives a correct value estimate but not, so far, an
exact reduction. If $m=\min f$ is attained at $p$ with $\norm p\le R$, then
$f_\varepsilon=f+\varepsilon\norm x^2$ is strongly convex and
$0\le\min f_\varepsilon-m\le\varepsilon R^2$. For a rational threshold $c$
with $m\ne c$, the sign of $\min f_\varepsilon-c$ equals that of $m-c$ once
$\varepsilon R^2<\abs{m-c}$. For this nonzero comparison gap, the general
separation bound, Lemma~\ref{lem:separation}(b) at degree four, guarantees
only $\abs{m-c}\ge2^{-2^{O(L)}}$. An $\varepsilon$ chosen from this uniform
bound can therefore require exponentially many printed bits; this does not
show that a particular instance needs such a small $\varepsilon$. Writing
$\varepsilon$ as a short circuit changes the input model of
Theorem~\ref{thm:exact-upper}, whose warm start depends on the bit length of
the curvature bound $2\varepsilon$. Formal infinitesimals do not repair the
iteration either. For $(x-1)^4+\varepsilon x^2$, Newton's method from $x_0=0$,
computed over $\Q(\varepsilon)$, has iterates with constant terms
$x_j(0)=1-(2/3)^j$, while the minimizer tends to $1$ as $\varepsilon\to0$; no
finite iterate is infinitesimally close to it. Equality tests also need a gap
shift, since regularization turns a zero minimum into a positive one.

These examples rule out three proof shortcuts. They do not show that exact
comparison for partially degenerate convex quartics lies outside the
complexity class of Theorem~\ref{thm:exact-upper}; that question is open.

\subsection{Facial restriction along a selected nullvector}
\label{app:bnd-nullvector}

Facial reduction simplifies a semidefinite feasibility problem by restricting
to matrices that vanish on a known kernel vector. In exact rational
computation one wants this step to preserve rational feasible points, for
instance when a rational Gram matrix is sought
(Sections~\ref{sec:heights}--\ref{sec:fields-certificates}). The safe
hypothesis is that the vector lies in the kernel of every feasible matrix. A
nullvector of one selected feasible matrix does not suffice.

\begin{proposition}[A selected nullvector can destroy rational feasibility]
\label{prop:bnd-nullvector}
Let
\[
\mathcal A(x)=\begin{pmatrix}2&x\\x&1\end{pmatrix}\oplus
\begin{pmatrix}x&1&0\\1&x&1\\0&1&x\end{pmatrix},\qquad
X(t,x)=\diag\bigl(t,\ \mathcal A(x)+tI_5\bigr).
\]
The rational affine family $X(t,x)$ contains the rational positive definite
matrix $X(2,0)$ and the positive semidefinite matrix $X(0,\sqrt2)$, which
annihilates $e_1$. Every positive semidefinite member that annihilates $e_1$
equals $X(0,\sqrt2)$, so the restricted family has no rational feasible
matrix.
\end{proposition}

\begin{proof}
The first block of $\mathcal A(x)$ is positive semidefinite exactly when
$\abs x\le\sqrt2$, and the second has eigenvalues $x-\sqrt2$, $x$ and
$x+\sqrt2$, so $\mathcal A(x)\succeq0$ exactly when $x=\sqrt2$. The matrix
$X(2,0)$ has eigenvalues $2,4,3,2-\sqrt2,2,2+\sqrt2$. A member that annihilates
$e_1$ has $t=0$, and then positive semidefiniteness forces $x=\sqrt2$.
\end{proof}

Here the restriction keeps the affine equations and adds $Xe_1=0$; deleting
the first row and column while leaving $t$ free is a different operation. The
selected matrix $X(0,\sqrt2)$ does not have maximal rank among feasible
matrices.

\begin{lemma}[Common kernels preserve rational feasibility]
\label{lem:bnd-common-kernel}
Let $\mathcal X$ be a set of positive semidefinite matrices defined by
rational affine equations. If a nonzero rational vector $\xi$ satisfies
$X\xi=0$ for every $X\in\mathcal X$, then a rational congruence maps
$\mathcal X$ bijectively onto a rational affine semidefinite family of smaller
size, preserving rational points in both directions. A nullvector of a
feasible matrix of maximal rank is annihilated by every feasible matrix.
\end{lemma}

\begin{proof}
Let $U$ be a rational matrix whose columns span $\xi^\perp$ and $C$ a rational
left inverse of $U$. Every symmetric $X$ with $X\xi=0$ satisfies
$X=UYU^{\mathsf T}$ with $Y=CXC^{\mathsf T}$, and both maps preserve
rationality and positive semidefiniteness. For the second statement, if
$G_0,G\succeq0$, then $\ker(G_0+G)=\ker G_0\cap\ker G$. If a feasible $G$ did
not annihilate some vector of $\ker G_0$, the midpoint of $G_0$ and $G$ would
be feasible with larger rank than $G_0$.
\end{proof}

Maximal rank does not by itself provide a rational basis of an irrational
common kernel, nor a way to find it. The safeguard is consistent with the
moment and Gram arguments of Section~\ref{sec:heights}, which identify the
common kernel through a feasible matrix of maximal rank.

\subsection{Singular equalities under uncertain data}
\label{app:bnd-uncertain}

The exact results of this paper read every coefficient exactly. Numerical
certification of equality-constrained problems often uses only enclosures of
function values. The two propositions below show that, at singular
equalities, this information can be insufficient in principle for an
accurate upper bound on the optimal value. Other conclusions can survive: in
the second example the existence of a root is robust under perturbation.
The propositions are statements about information, not lower bounds for
algorithms that receive exact rational input: exact symbolic data remove the
uncertainty in both examples.

On a compact set $\mathcal K$, let
$v(\mathfrak h)=\min\{\phi(x):x\in\mathcal K,\ \mathfrak h(x)=0\}$ for an
objective $\phi$ and an equality map $\mathfrak h$ with nonempty zero set. If
the available information allows every map in a family $\mathfrak H$, then an
upper bound justified by this information alone must be at least
$\sup_{\mathfrak g\in\mathfrak H}v(\mathfrak g)$. Each permitted map needs a
good feasible point of its own; no common point is required.

\begin{proposition}[A gap that does not close]
\label{prop:bnd-fixed-gap}
Let $\mathcal K=[-1,2]$, $\phi(x)=x$, $\mathfrak h(x)=x^2(x-1)$, and let
$\mathfrak H_\delta$ consist of the continuous functions within $\delta$ of
$\mathfrak h$ in the maximum norm, $0<\delta<2$. Let $r_\delta\in(1,2)$ solve
$r_\delta^2(r_\delta-1)=\delta$. Then every member of $\mathfrak H_\delta$ has a
zero in $\mathcal K$, $v(\mathfrak h)=0$, $\sup_{\mathfrak g\in\mathfrak
H_\delta}v(\mathfrak g)=r_\delta$, and $r_\delta\to1$ as $\delta\to0$.
\end{proposition}

\begin{proof}
On $[1,2]$, $\mathfrak h$ increases from $0$ to $4$, so $r_\delta$ exists and
tends to $1$. Every permitted $\mathfrak g$ has $\mathfrak g(-1)\le-2+\delta<0$
and $\mathfrak g(r_\delta)\ge0$, so it has a zero in $[-1,r_\delta]$. The
member $\mathfrak h-\delta$ is negative on $[-1,1]$ and has the single zero
$r_\delta$ in $(1,2]$.
\end{proof}

No uniform enclosure of positive width therefore justifies an upper bound
below one, although the true optimum is zero. The optimum is a zero of even
multiplicity, and a nearby simple zero of a perturbation does not converge to
it. Exact factorization certifies $x=0$ at once.

The second example keeps the singular root robust and still needs exponential
precision.

\begin{proposition}[Exponential precision at a robust root]
\label{prop:bnd-robust-root}
For $k\ge1$, on $[-1,1]^{k+1}$ with variables $(y_1,\ldots,y_k,x)$ and
$y_0=x$, consider
\[
v_k(a)=\min\{x^2:\ y_i-y_{i-1}^2=0\ (1\le i\le k),\ xy_k=a\},\qquad
\mathcal D_k=2^k+1 .
\]
\begin{enumerate}[label=(\alph*),leftmargin=*]
\item For $\abs a\le1$ the feasible set is the single point
$x=\operatorname{sgn}(a)\abs a^{1/\mathcal D_k}$,
$y_i=\abs a^{2^i/\mathcal D_k}$, and $v_k(a)=\abs a^{2/\mathcal D_k}$.
\item At $a=0$ the root is isolated, the Jacobian has rank $k$, and the local
Brouwer degree is $+1$; every sufficiently small continuous perturbation of
all equations has a zero near the origin.
\item Every equation is quadratic in at most two variables with coefficients
in $\{-1,0,1\}$ at $a=0$, and the interaction graph has treewidth at most two.
\item If the last constant is known only to lie in $[-\delta,\delta]$,
$0\le\delta\le1$, and the other equations are exact, the least valid upper
bound is $\delta^{2/\mathcal D_k}$. An upper bound within $\varepsilon\in(0,1)$ of
the true value $v_k(0)=0$ therefore requires $\delta\le
\varepsilon^{\mathcal D_k/2}$, that is, $\log_2(1/\delta)\ge\tfrac{2^k+1}2
\log_2(1/\varepsilon)$.
\end{enumerate}
\end{proposition}

\begin{proof}
The chain gives $y_i=x^{2^i}$ and $x^{\mathcal D_k}=a$; an odd power is a
bijection of $\R$, which gives (a). At the origin the first $k$ rows of the
Jacobian contain the identity in the $y$ columns and the last row is zero.
For a small target $(0,\ldots,0,a)$, $a\ne0$, the only preimage is the point
in (a), and eliminating the triangular $y$ block shows that the Jacobian
determinant there is $\mathcal D_kx^{\mathcal D_k-1}>0$. The degree is therefore
$+1$ by the regular-value formula, and homotopy invariance of the degree gives
the robustness in (b)~\cite{Deimling1985}. The interaction graph is the path
$x,y_1,\ldots,y_k$ with the closing edge $y_kx$. For $k\ge2$, the bags
$\{x,y_i,y_{i+1}\}$, $1\le i<k$, arranged along a path, form a tree
decomposition of width two; for $k=1$ the single bag $\{x,y_1\}$ has width
one. Part (d) follows
from (a), since $\abs a^{2/\mathcal D_k}$ is largest at $\abs a=\delta$.
\end{proof}

All derivatives of positive order are independent of $a$, so exact
derivative data do not help; the exact statement $a=0$ does, since the origin
is then an immediate certificate. The robustness in (b) concerns existence,
not uniqueness: replacing the first equation by $y_1-x^2=-\eta$ produces three
solutions near the origin. The bound in (d) is a statement about an absolute
uncertainty radius, not about the running time of any algorithm, and it
grows with the chain length rather than with the binary length of exact
input. Exponential error-bound exponents for such power chains are
classical~\cite{Kollar1999}; robustness of zeros under perturbation is
studied in~\cite{FranekRatschanZgliczynski2016}.

\subsection{A succinct nonconvex root penalty}
\label{app:bnd-penalty}

Exact penalties replace constraints by a term in the objective. For bounded
quadratic problems, a fractional power of the constraint violation gives an
exact penalty whose coefficient and exponent both have polynomial binary
length, and the fractional power has a polynomial-size quadratic lift. The
lift is nonconvex. Eliminating its squaring chain leaves polynomial
inequalities of exponential degree, and eliminating all auxiliary variables
returns the fractional power, which is not a polynomial. None of these forms
satisfies the hypotheses of Theorem~\ref{thm:exact-upper}. The result marks
the difference between the input models of this paper: explicit low-degree
expansions, high-degree eliminated forms, and circuits.

The setting is explicitly expanded quadratic input. Let $X\subset\R^N$ be
compact, described by rational quadratic equations, weak quadratic
inequalities, finite rational variable bounds, and integrality of some
coordinates. Let $\phi$ be a rational quadratic objective, and let further
rational quadratic equations and weak inequalities define a nonempty
$S\subseteq X$. All polynomials are given by their coefficients, and $L$ is
the total binary length. The \emph{signed residuals} are $h(x)$ for each
additional inequality $h(x)\le0$ and $\pm h(x)$ for each additional equation
$h(x)=0$. Let $V(x)$ be the maximum of zero and the signed residuals, that
is, of zero, the positive parts of the additional inequalities, and the
absolute values of the additional equations. Termwise estimates on the
bounding box give positive
integers $H_V,B_\phi$ of polynomial bit length with $V\le H_V$ and
$\abs\phi\le B_\phi$ on $X$; put $\mathcal R=V/H_V\in[0,1]$.

Two effective theorems of semialgebraic geometry are used, with the bounds
of Basu and Mohammad-Nezhad. As in the source, the degree bound $\mathfrak d$
is at least two; a description of lower degree is used with $\mathfrak d=2$.
First~\cite[Theorem~4.1]{BasuMohammadNezhad2024}, a formula with one block of
$\kappa$ quantified variables, $\ell$ free variables, and polynomials of
degree at most $\mathfrak d$ with integer coefficients of bit length at most
$\mathfrak t$ is equivalent to a quantifier-free formula whose polynomials
have degree $\mathfrak d^{O(\kappa)}$ and integer coefficients of bit length
$\mathfrak t\,\mathfrak d^{O(\kappa)O(\ell)}$.
Second~\cite[Theorem~2.2]{BasuMohammadNezhad2024}, if
$\mathcal A\subset\R^{\mathfrak n}$ is compact and described with degrees at
most $\mathfrak d$ and integer coefficients of bit length at most $\mathfrak t$,
and $\mathfrak g,\mathfrak e\colon\mathcal A\to[0,\infty)$ are continuous with
graphs of the same kind and $\mathfrak e^{-1}(0)\subseteq\mathfrak g^{-1}(0)$,
then $\mathfrak g^{\varrho}\le C\mathfrak e$ on $\mathcal A$ for a positive
integer $\varrho\le(8\mathfrak d)^{2(\mathfrak n+7)}$ and a constant $C>0$
with $\log_2\max(1,C)\le\mathfrak t\,\mathfrak d^{O(\mathfrak n^2)}$.

\begin{theorem}[Succinct root penalty]
\label{thm:bnd-root-penalty}
There is a polynomial $\mathfrak p$ with nonnegative integer coefficients and
$\mathfrak p(0)\ge1$, independent of the instance, such that, with
$\varpi=2^{-\mathfrak p(L)}$,
\[
\min_{x\in X}\bigl\{\phi(x)+(4B_\phi+1)\mathcal R(x)^{\varpi}\bigr\}
=\min_{x\in S}\phi(x),
\]
and the two problems have the same minimizers. Moreover, with variables
$s_0,\ldots,s_{\mathfrak p(L)}\in[0,1]$, the problem of minimizing
$\phi(x)+(4B_\phi+1)s_0$ over $x\in X$ subject to
\[
s_{j+1}=s_j^2\quad(0\le j<\mathfrak p(L)),\qquad
H_Vs_{\mathfrak p(L)}\ge V(x),
\]
where the last condition is one quadratic inequality for each signed
residual, is an exact quadratic lift of the penalized problem with
polynomially many variables and coefficients of polynomial bit length.
\end{theorem}

\begin{proof}
For an integer coordinate $z$ with rational bounds $l\le z\le u$, put
$\ell_z=\lceil l\rceil$ and $u_z=\lfloor u\rfloor$; the integers in $[l,u]$
are those in $[\ell_z,u_z]$, and $\ell_z\le u_z$ because $S$ is nonempty.
Write $z=\ell_z+\sum_{j=0}^{J_z}2^jb_j$ with $b_j(1-b_j)=0$ and
$0\le b_j\le1$, where $J_z=\lfloor\log_2(u_z-\ell_z)\rfloor$ (no bits if
$u_z=\ell_z$), and keep the row $z\le u_z$. This gives a compact real
quadratic lift $X'$ of polynomial size whose projection is $X$; let $S'$ be
the lift of $S$. Let
$\vartheta=\min_S\phi\in[-B_\phi,B_\phi]$, attained by compactness, and
\[
\mathcal I=\{t\in[-B_\phi,B_\phi]:\ \text{no }y\in S'\text{ has }\phi(y)<t\}
=[-B_\phi,\vartheta].
\]
This set is described, after clearing denominators, with one quantifier
block of polynomially many variables, degree two, and coefficients of
polynomial length, so the first effective theorem, with $\mathfrak d=2$ and
$\ell=1$, describes it without quantifiers with degrees and coefficient
lengths $2^{\poly(L)}$. The unknown, possibly irrational, value $\vartheta$ is
never used as a coefficient.

On $\mathcal A=X'\times\mathcal I$ let $\mathfrak g(w,t)=\max\{0,t-\phi(w)\}$
and $\mathfrak e(w,t)=\mathcal R(w)$. Their graphs are described by the
conditions defining $\mathcal A$ together with Boolean combinations of
quadratic conditions, so they obey the same degree and coefficient bounds.
If $\mathfrak e(w,t)=0$, then $w\in S'$ and $t\le\vartheta\le\phi(w)$, so
$\mathfrak g(w,t)=0$. The second effective theorem gives
$\mathfrak g^\varrho\le C\mathfrak e$ with $\varrho,\log_2\max(1,C)\le
2^{\mathfrak a(L)}$ for a polynomial $\mathfrak a$ independent of the
instance. Choose $\mathfrak p$ with nonnegative integer coefficients,
$\mathfrak p(0)\ge1$ and $\mathfrak p\ge\mathfrak a$ on $[1,\infty)$; then
$\varpi\varrho\le1$ and $\varpi\log_2\max(1,C)\le1$. Since
$\mathfrak g\le2B_\phi$,
\[
\mathfrak g\le\min\{2B_\phi,\,C^{1/\varrho}\mathfrak e^{1/\varrho}\}
\le(2B_\phi)^{1-\varpi\varrho}C^{\varpi}\mathfrak e^{\varpi}
\le4B_\phi\,\mathfrak e^{\varpi},
\]
and the bound also holds where $\mathfrak e=0$. At $t=\vartheta$ this reads
$\phi(x)+(4B_\phi+1)\mathcal R(x)^{\varpi}\ge\vartheta+\mathcal
R(x)^{\varpi}$ on $X$. Infeasible points are strictly worse than $\vartheta$,
and feasible minimizers attain it.

In the lift, nonnegativity and the squaring equations give
$s_{\mathfrak p(L)}=s_0^{2^{\mathfrak p(L)}}$, so for fixed $x$ the least
feasible $s_0$ is $\mathcal R(x)^{\varpi}$, and the positive coefficient of
$s_0$ selects it.
\end{proof}

Exact penalties with fractional powers of the violation on compact
semialgebraic sets are classical~\cite{JiaoPhamTuyen2025}; the theorem adds
the encoding bound for explicitly expanded quadratic input and the quadratic
lift. The constants are effective but have not been evaluated. The theorem is
a representation statement and gives no efficient optimization method. Its
scope is limited in four ways.

\begin{proposition}[Limits of the root penalty]
\label{prop:bnd-penalty-limits}
Let $k\ge1$ be the chain length in (b)--(d).
\begin{enumerate}[label=(\alph*),leftmargin=*]
\item The squaring equations cannot be relaxed to $s_{j+1}\ge s_j^2$: then
$s_0=\cdots=s_{\mathfrak p(L)-1}=0$ and $s_{\mathfrak p(L)}=\mathcal R(x)$ is
feasible and incurs no penalty. Eliminating the chain instead gives the
polynomial inequalities $H_Vs_0^{2^{\mathfrak p(L)}}\ge h(x)$, one for each
signed residual $h$, of exponential degree.
\item On the native curve $x_{i+1}=x_i^2$, $0\le x_i\le1$ ($0\le i<k$), with
objective $-x_0$ and additional equation $x_k=0$, a penalty
$\rho\abs{x_k}^{\varpi}$ is exact only if $\varpi\le2^{-k}$, for every
finite $\rho$.
\item With native $z\in\{0,1\}$, $y_0=\tfrac12$, $y_{i+1}=y_i^2$, all
variables in $[0,1]$, objective $-z$ and additional equation $zy_k=0$, the
penalty $\rho\abs{zy_k}^{\varpi}$ has the same minimizers exactly when
$\rho>2^{\varpi2^k}$.
\item In (b), at $x_0=\tfrac12$ the residual is $r=2^{-2^k}$. A lift of depth
at least $k$ with all auxiliary variables zero except the last, which equals
$r$, violates only the last squaring equation, by $r$. A feasibility tolerance
of at least $r$ therefore accepts zero penalty and objective $-\tfrac12$,
although the exact minimum is zero.
\end{enumerate}
\end{proposition}

\begin{proof}
(a) is immediate. In (b) the penalized objective along the curve is
$-\sigma+\rho\sigma^{\varpi2^k}$ with $\sigma=x_0$, which is negative for
small $\sigma>0$ when $\varpi2^k>1$. In (c) the two native points are
$z=0$, with value $0$, and $z=1$, with value $-1+\rho2^{-\varpi2^k}$. Part (d)
is a direct substitution.
\end{proof}

The theorem requires an explicit quadratic objective with a bound $B_\phi$ of
polynomial length. For the sparse objective $x^{2^k}$ on $[1,2]$ with the
additional constraint $x\ge2$, the violation $[2-x]_+$ equals one at $x=1$
for every exponent, while the objective gap is $2^{2^k}-1$; no exponent
avoids a coefficient of that size. This example lies outside the theorem's
input model; it does not show that the theorem accepts high binary degree.

The lift has quadratic equations of the form $s_{j+1}=s_j^2$, so it is not a
native PSD system. Eliminating $s_1,\ldots,s_{\mathfrak p(L)}$ while keeping
$s_0$ leaves the polynomial inequalities of part~(a), of degree
$2^{\mathfrak p(L)}$; eliminating $s_0$ as well returns the term
$\mathcal R^{\varpi}$, which is not a polynomial. None of these forms is the
unconstrained minimization of an explicitly given globally convex polynomial
of bounded degree, so Theorem~\ref{thm:exact-upper}, which needs such an
objective with a strong-convexity bound or certificate, does not apply.
Part (c) with $\varpi=1$
shows that an exponent-one penalty can need a coefficient with exponentially
many bits in the chain length; the theorem trades such coefficients for a
small exponent and nonconvex auxiliary equations.
